\section{Analytic entry and the activity coordinates}\label{sec:entry}

We now connect the finite height laws to the small-activity map of
\cite{F}.  The map itself acts on analytic functions and is independent
of the temperature of the physical input.  The point requiring a new
argument is initialization: the critical construction must remain small
on a parameter disk large enough to differentiate with respect to $b$.
We prove this on a disk of radius proportional to $s^{-2}$, where $s$ is
the observation-cell side.  All constants depending on fixed geometric
parameters are chosen before $s$ tends to infinity.

\subsection{The reference field and the initialization norm}

The critical free-height coefficient is $a=8\pi$ by
\cite[Section~6]{X}.  Fix $\alpha=\sqrt a$ and $\omega=2\pi/\alpha$.
The reference field has covariance $A^+$, so that $\alpha$ times the
reference field has height covariance $aA^+$.  We keep this reference,
including $\alpha$, fixed when $b$ varies.

Partition the primary graph into whole square cells of side $s$, let
$B_s$ average over each cell, and introduce the observations
\begin{equation}\label{eq:observations}
                         \eta=B_sh+e,
\end{equation}
where the components of $e$ are independent standard Gaussians.  On a
retained graph $G$, let $Z_{G,b}(\eta)$ be the height partition sum with
penalty $\exp(-\|B_sh-\eta\|^2/2)$, and define the fixed reference energy
\begin{equation}\label{eq:observation-energy}
 E_G(\eta)=\min_{v\in\R^{V(G)}}
       \{a^{-1}\|\nabla v\|_G^2+\|B_sv-\eta\|^2\}.
\end{equation}
Each component constant is free in this minimization.  For a retained
cell face $f=(Q,Q')$, put $w_f=\eta_{Q'}-\eta_Q$ and use the entire
cutoff
\begin{equation}\label{eq:entry-cutoff}
 \chi_s(z)=
 \frac{\int_{-s}^s e^{-(z-t)^2/2}\,dt}{\int_{-s}^se^{-t^2/2}\,dt},
 \qquad
 \widehat Z_{G,b}(\eta)=Z_{G,b}(\eta)\prod_f\chi_s(w_f).
\end{equation}
A deleted face loses both its primary bonds and its cutoff factor.

Here are the reference split and norm used in \cite{F}.  For an integer
$u\ge1$, set
\[
 P_u(\cos t)=\left(\frac{\sin(ut/2)}{u\sin(t/2)}\right)^{64},
 \qquad \mathsf P_u=P_u(1-A/64),
\]
with removable values understood, and put
\begin{equation}\label{eq:reference-split}
 C_{<u}=A^{-1}(I-\mathsf P_u),\qquad
 \Gamma_u=A^{-1}(\mathsf P_u-\mathsf P_{Lu}),\qquad
 C_{\ge u}=A^+\mathsf P_u.
\end{equation}
The first two quotients have polynomial continuations at the constant
mode.  They have finite range and are positive.  Their constant modes
can be integrated because the final common mean is uniform modulo
$\omega$.

At side $u$, a polymer is a finite set $X$ of $u$-blocks connected for
the fixed adjacency of \cite{F}.  Its occupied support is the union of
those blocks.  The field dependence of $K(X,\phi)$ is confined to this
support together with the prescribed microscopic collar of
\cite[Proposition~4.1]{F}, whose width is at most a fixed sufficiently
small fraction of $u$.  This dependence region is distinct from the
larger norm padding $D_u(X)$, which is enlarged to whole observation
cells and contains all fixed difference stencils.

We use the bulk symmetry class of \cite[Theorem~4.6(i)]{F}.
Its activities are analytic in the field and satisfy
\begin{equation}\label{eq:bulk-activity-symmetries}
 K(X,-\phi)=K(X,\phi),\qquad
 K(X,\phi+\omega\one)=K(X,\phi),\qquad
 K(gX,g\phi)=K(X,\phi),
\end{equation}
where $g$ is a compatible block translation or square-lattice rotation
or reflection and $(g\phi)(v)=\phi(g^{-1}v)$.  Compatibility means
that $g$ preserves the chosen block and observation grids.  On a torus
we impose the corresponding periodic support and covariance rules.
These identities are complex-linear conditions on the activity;
reality on real fields is a separate property.  The point norm is
\begin{equation}\label{eq:entry-point-norm}
 \begin{split}
 \|f\|_{u,D}&=\max_{v\in D}\{|f(v)|,u|\nabla f(v)|,
                                           u^2|\nabla^2f(v)|\},\\
 |K(X)|_{u,\phi}
 &=\sum_{r\ge0}\frac{h_0^r}{r!}
   \sup_{\|f_i\|_{u,D_u(X)}\le1}
       |D^rK(X,\phi)[f_1,\ldots,f_r]|.
 \end{split}
\end{equation}
Let $\mathscr W_u^\kappa(X,\phi)$ be the ordinary exponential regulator
specified in \cite[Section~4.2, Equations~(41)--(43)]{F}; it controls the bulk gradient energy,
its boundary trace, and scaled second differences.  The additional
observation regulator is given explicitly by
\begin{equation}\label{eq:entry-observation-regulator}
 V_u(X,\phi)=\frac12\sup_{\zeta,e}
 \left\{h_*E_{D_u(X)}\bigl(B_s\alpha(\phi+\zeta)+e\bigr)
       -p_1^{-1}(\|\zeta\|_{C_{<u}}^2+\|e\|^2)\right\}.
\end{equation}
The Cameron norm is taken on the covariance support.  With a fixed block
weight $A_w>1$, the activity norm is
\begin{equation}\label{eq:entry-activity-norm}
 \|K\|_{u,s}=\sup_{X,\phi}
 A_w^{|X|}\frac{|K(X)|_{u,\phi}}
                  {\mathscr W_u^\kappa(X,\phi)e^{V_u(X,\phi)}}.
\end{equation}
Write $\mathscr A^{\mathrm{bulk}}_{u,s}$ for the complex space of
activity families satisfying the support and symmetry conditions above
and having finite norm \eqref{eq:entry-activity-norm}.  Every activity
ball and activity variation below is taken in this space or its
compatible-torus counterpart.  The ordinary regulator is kept exactly
as in the cited construction; its shell averaging and translation estimates, together with the
explicit observation regulator, are part of the imported local-map
bounds below.  Norms with different observation sides $s$ need not be
identified.

The cutoff gives a strict large-field reserve.  Indeed, adjacent-cell
averages have a primary flow representation with coefficient bound
$C/s$ and squared cost $C$.  Applying
\eqref{eq:bounded-flow-mgf} with tilts of size at most a fixed multiple
of $s$, then using the Gaussian tails of $\chi_s$, gives uniformly for
$b$ in a compact neighborhood of $\bc$
\begin{equation}\label{eq:entry-large-field}
 e^{E_G(\eta)/2}\frac{\widehat Z_{G,b}(\eta)}{Z_{G,b}(0)}
             \le e^{(1-c_*)E_G(\eta)/2},\qquad \eta\in\R^{\mathcal Q(G)},
\end{equation}
for some fixed $c_*>0$ on every retained topology used in the
construction.  More explicitly, transporting each site straight to its
translate in a neighboring cell gives flows $\Phi_f$ with
\[
 \max_e\left|\sum_ft_f\Phi_f(e)\right|\le Cs^{-1}\max_f|t_f|,
 \qquad \sum_e\left|\sum_ft_f\Phi_f(e)\right|^2\le C\sum_ft_f^2.
\]
The tilted affine comparison applies also at prescribed observations;
releasing the height and noise equalities gives a bound
$\exp(C\sum_ft_f^2)$.  Choose
$t_f$ proportional to $\operatorname{sgn}(w_f)\min(|w_f|,s)$.
For $|w_f|>s$ the cutoff supplies the remaining quadratic decay.
Finally $E_G(\eta)\le C\sum_fw_f^2$, by cell interpolation.
These are the steps in \cite[Lemma~4.4 and Equation~(57)]{F}; the only physical transform
used there is $\exp\{b(\cosh(2\pi t)-1)\}$, whose constants are uniform
on the chosen $b$ interval.

Choose fixed parameters with strict inequalities
\begin{equation}\label{eq:entry-reserves}
        1-c_*<h_1<h_*<1/p_2,\qquad 1<p_1<p_2.
\end{equation}
There is also slack in the exponent of $\mathscr W$.  Choose the fixed
analytic radius, field-strip widths, and a sufficiently large dyadic
$L$ as in \cite{F}; then choose the activity weights and the small
map domain.  Fix its coloring period $q$, choose a sufficiently large
dyadic $R_0$, and set $P_0=R_0^5$.  An admissible observation side
satisfies
\begin{equation}\label{eq:entry-side}
                         qR_0s=L^{j_0}.
\end{equation}
The index $j_0$ is absolute, rather than reset to zero at entry.
The choice of $R_0$ remains subject to any further fixed small-tube
requirements made below; all such choices precede the limits in $s$
and $b$.

\subsection{A parameter disk for initialization}

We use the finite initialization of \cite{F} with no magnetic or
compulsory component.  On a compatible finite torus it represents the
cutoff law exactly, up to a field-independent scalar, after averaging
$e$ and the Gaussian field with covariance $C_{<L^{j_0}}$.  The
corresponding plane activities are defined by the same local formulas.
For clarity, the partition functional of an activity is
\[
 \mathcal Z_u(K;\phi)=\sum_{V\subset\mathcal B_u}
                         \prod_{X\in\operatorname{Comp}(V)}K(X,\phi),
\]
where components use the fixed block adjacency and the empty product
is one.  The representation identifies the averaged cutoff interaction
with a scalar times $\mathcal Z_{L^{j_0}}$.

\begin{proposition}[Analytic initialization]\label{prop:entry}
For every sufficiently small $\rho>0$, the fixed choices above can be
made so that there exist $c_{R_0}>0$ and $s_0<\infty$ with the following
property.  For every admissible $s\ge s_0$, the ordinary plane and
compatible-torus entry activities $K_s(b)$ are Banach analytic on
\begin{equation}\label{eq:entry-parameter-disk}
                         |b-\bc|<c_{R_0}s^{-2},
\end{equation}
and take values in $\mathscr A^{\mathrm{bulk}}_{L^{j_0},s}$ or its
compatible-torus counterpart, with norm at most $\rho$.
They are real on real fields when $b$ is real.  Their support,
symmetry, and local copying identities agree with those at criticality
throughout the complex parameter disk.
For real parameters in this disk the cutoff partition representation
is exact.  On a smaller concentric disk, for each fixed $r\ge1$,
\begin{equation}\label{eq:entry-cauchy}
                \|\partial_b^rK_s(b)\|_{L^{j_0},s}
                              \le C_{R_0,r}s^{2r}.
\end{equation}
\end{proposition}

\begin{proof}
We first describe the parts of the finite expansion whose estimates
have to be extended.  The geometry and algebra are those of
\cite[Sections~4.6--4.7, Theorem~4.6]{F}; only their physical weights change with $b$.
The observation cells are partitioned into tiles of side comparable to
$R_0$.  A finite colored sequence of isolation events deletes the
remaining bonds between tiles.  At an event with open patch $o_i$ and
cut patch $c_i$, its proxy is
\begin{equation}\label{eq:entry-proxy}
 \mathcal B_i(b,\eta)=C_i(b)
                  e^{-(E_{o_i}(\eta)-E_{c_i}(\eta))/2},\qquad
 C_i(b)=\frac{Z_{o_i,b}(0)}{Z_{c_i,b}(0)}.
\end{equation}
One writes the original factor as its proxied cut factor plus their
difference.  Selecting the difference marks that event.  Later events
within its fixed neighborhood are skipped.  After the finite sequence,
undeleted intertile bonds lie in the union of the marked neighborhoods.
We call each connected part of this union a halo.  Outside the halos,
an isolated tile contributes either one or its normalized factor minus
one.  The latter choice is a selected tile.

The remaining Gaussian mismatch is expanded in links
$e^{b_\ell w_fw_{f'}}-1$, with
\begin{equation}\label{eq:entry-links}
                         |b_\ell|\le Ce^{-c(R_0+d_\ell)}.
\end{equation}
Here $d_\ell$ records the entire tagged path length, including winding
on a torus.  A link covers at most $C(1+d_\ell/R_0)$ entry blocks.
An individual pattern $\pi$ consists of marked events, selected tiles,
and links, with their complete occupied cover $X(\pi)$.  Write $l(\pi)$
for the number of marks and selected tiles.  Each of these physical
items uses at most $C_{R_0}$ observation cells and $C_{R_0}s^2$ primary
edges.  Link-only cells carry no physical height factor.  All identical
factors off the affected halos cancel before the pattern is defined.
Its denominator is a product of the affected local scalars $C_i(b)$
and isolated-tile normalizers $Z_{Q,b}(0)$.  In particular, a centered
partition function of an arbitrarily large halo is not a denominator
of a defining pattern.

We keep every Gaussian energy in \eqref{eq:entry-proxy} fixed.  Thus the
finite algebra, the link coefficients, the occupied covers, and the
copying identities remain valid with the actual parameter in every
physical factor and scalar normalizer.  It remains to bound and sum the
resulting patterns.

\paragraph{Real parameters.}
In the critical proof a positive expansion of each marked difference
introduces a conditional factor
\begin{equation}\label{eq:entry-mark-fraction}
 d_i(\eta,\gamma)=
 \frac{|K_{o_i}(\eta,\gamma)-\mathcal B_i(\eta)K_{c_i}(\eta,\gamma)|}
      {K_{o_i}(\eta,\gamma)+\mathcal B_i(\eta)K_{c_i}(\eta,\gamma)}
                         \le1.
\end{equation}
The kernels prescribe the primary guard values $\gamma$ surrounding
the event.  Exterior factors common to the two choices cancel in this
ratio.  The halo is then cut into boxes of side $P_0$ in cell units.
Only guards whose testing neighborhoods lie inside a box are retained.
Affine comparison bounds the joint observation and guard densities
before and after cutting.  The resulting centered zero-guard ratio is
bounded by a product of local separated-pin comparisons.  Gaussian
shape errors have small face-square coefficients.  This is precisely
the positive-topology and seam accounting in
\cite[Equations~(71)--(74)]{F}.

After $R_0,P_0$, the field strips, and the large-field threshold are
fixed, the compact-data part of this argument concerns finitely many
box topologies.  Each has at most $C_{R_0}s^2$ primary edges.  A box is
called good when the sum of squared differences on its underlying cell
faces is below that fixed threshold.  Period recentering places its
observations in a fixed compact set.  The critical inputs in
\cite[Proposition~2.3 and Equations~(75)--(80)]{F} give on these compact sets: bounds on normalized positive
box factors, separated-pin restoration bounds, and a small fixed list
of moments of the fractions \eqref{eq:entry-mark-fraction}.  The list
includes all H\"older orders required by the finite number of marks
which a $P_0$-box can contain.  All inequalities are used with strict
slack.

Lemma~\ref{lem:finite-perturbation}, applied to these finite graphs,
preserves the required tolerances when
$|b-\bc|\le c_{R_0}s^{-2}$.  For centered pin interactions one applies
the lemma to each pinned sum itself, so that a small absolute change
in an unpinned normalization is never used as a relative estimate for
a rare pin event.  On real compact observation sets,
\[
 Z_{G,b}(\eta)/Z_{G,b}(0)
 =e^{-\|\eta\|^2/2}\E_{G,b,0}e^{\langle\eta,B_sh\rangle}
                       \ge e^{-\|\eta\|^2/2},
\]
which converts normalized total-variation estimates to the needed
relative estimates.

The conditional fractions require some care because the guards have
an increasing number of primary vertices.  For one retained event,
use the two complete adjacent positive box topologies and let
$A_b(\gamma),B_b(\gamma)$ be their joint weights, incorporating the
proxy and cutoff scalars.  Their total masses are comparable, uniformly
on the fixed compact set: at zero observation this follows from the
open and zero-guard brackets and the bounded separated-pin interaction;
for other compact observations it follows from the preceding lower
bound and the upper affine comparison.  Lemma~\ref{lem:finite-perturbation}
controls both arrays in $\ell^1$, in units of either total mass.
For nonnegative $A,B,A',B'$,
\[
 (A+B)\left|
 \frac{|A-B|}{A+B}-\frac{|A'-B'|}{A'+B'}\right|
 \le2(|A-A'|+|B-B'|),
\]
with the continuous convention at a zero denominator.  The same
estimate for the $r$th power has an additional factor at most $r$.
Consequently every required fixed moment of $d_i$ changes by a quantity
tending to zero with $c_{R_0}$.  No estimate uniform in individual guard
values is needed.  The compact imaginary-observation estimates after
pinning strips are preserved by the complex version of
Lemma~\ref{lem:finite-perturbation}.

The remaining inputs are uniform finite comparisons and
\eqref{eq:entry-large-field}.  They control the boxes outside the
compact set with the same reserves.  The deterministic Gaussian
localization, link estimates, and counting do not involve $b$.
Repeating the real and complex-field interpolation of
\cite[Lemma~4.7, Equations~(85)--(89)]{F} therefore gives, on fixed field disks with radius
larger than the polarization radius needed in
\eqref{eq:entry-point-norm},
\begin{equation}\label{eq:entry-real-pattern}
 |F_\pi(b;\phi+zf)|
 \le C^{l(\pi)}e^{-cR_0l(\pi)}
       \prod_{\ell\in\pi}Ce^{-c(R_0+d_\ell)}
       \mathscr W_{L^{j_0}}^\kappa(X,\phi)e^{V_{L^{j_0}}(X,\phi)}.
\end{equation}
Here $b$ is real in the parameter interval, $\phi$ is real,
$\|f\|_{L^{j_0},D}\le1$, and $|z|$ is in that fixed disk.
The constants in the exponential gains may have decreased.

We have extended the small real-parameter bound.  To obtain parameter
Cauchy estimates we next need an analytic bound on a complex disk; this
bound need not retain the small factor attached to each physical item.

\paragraph{A complex-parameter majorant in the same norm.}
Let $r_s=c_{R_0}s^{-2}$ and $b_+=\bc+r_s$.  For $|b-\bc|\le r_s$ we
have $|b|\le b_+$.  On a positive physical branch containing $N_\pi$
retained primary edges, absolute augmented Poisson weights satisfy
\begin{equation}\label{eq:complex-poisson-majorant}
 \left|\prod_e e^{-b}\frac{(b/2)^{U_e+V_e}}{U_e!V_e!}\right|
 \le e^{(b_+-\Re b)N_\pi}
       \prod_e e^{-b_+}\frac{(b_+/2)^{U_e+V_e}}{U_e!V_e!}.
\end{equation}
The physical support count gives $N_\pi\le C_{R_0}l(\pi)s^2$.
Each local scalar normalizer is nonzero in this disk by
Lemma~\ref{lem:finite-perturbation}; its value and its inverse, relative
to the corresponding value at $b_+$, cost at most a fixed factor.
There are $O(l(\pi))$ affected scalars.  Expand each marked difference
into its two terms and compare their unnormalized physical numerators
by \eqref{eq:complex-poisson-majorant}, and their local scalar factors
separately.  The result is bounded by $e^{C_{R_0}l(\pi)}$ times the
corresponding positive pattern at $b_+$.  At this point one may apply
the real positive-topology estimate at $b_+$.  This ordering uses no
complex partition function of a large connected halo.

For a field displacement $zf$ in the fixed field disk, every imaginary
observation coordinate is bounded.  The absolute observation penalty
costs $e^{\|\Im\eta\|^2/2}$ on each physical factor.  There are at
most $C_{R_0}l(\pi)$ such coordinates, including the local proxies.
The real quadratic forms satisfy
\[
                 \Re E_G(\eta_1+i\eta_2)=E_G(\eta_1)-E_G(\eta_2).
\]
Their imaginary contributions therefore also cost at most
$e^{C_{R_0}l(\pi)}$.  For the real parts keep the Gaussian shapes
combined into the reference mismatch before estimating them.  It
consists of the energy of the actual positive physical topology and
the small localized proxy and link errors.  Thus overlapping proxy
patches do not create overlapping broad-energy charges.

More explicitly, after dropping the fractions $d_i$, the large-field
bound may be used on every actual halo piece and every separate
selected-tile factor.  Its broad quadratic cost is
\begin{equation}\label{eq:complex-broad-charge}
                   \frac{1-c_*}{2}\sum_B E_B(\eta_1).
\end{equation}
The actual graphs $B$ are disjoint: selected tiles lie outside the halo
graph and are mutually isolated.  Restricting a trial field for the
open energy on $D=D_{L^{j_0}}(X)$ to each of them gives
$\sum_BE_B\le E_D$.  The remaining proxy, seam, and link costs have
small coefficients on underlying face differences with bounded
per-face sum.  The cell flow estimate bounds each such face square
by the energy of its two underlying cells.  Taking $R_0$ large leaves
these costs inside the strict gap $h_1-(1-c_*)$.  Hence the complete
positive real-variable quadratic form $q_\pi$ obeys
\begin{equation}\label{eq:complex-form-domination}
                           q_\pi(\eta_1)\le\tfrac12h_1E_D(\eta_1).
\end{equation}
This is the form domination of \cite[Equation~(82)]{F}, with all physical boxes
now charged broadly.  Its coefficient and its ambient open energy are
unchanged.

To see that the determinant is still harmless, write the averaging
variables as a standard Gaussian vector $\zeta$ and
$\eta_1=d+\mathsf S\zeta$.  Set
$M_\pi=\nabla_\zeta^2[p_1q_\pi(d+\mathsf S\zeta)]$.
The covariance contraction for the observation energy and
\eqref{eq:complex-form-domination} give
$0\le M_\pi\le p_1h_1I<I$.  Gaussian completion is exactly
\begin{equation}\label{eq:complex-entry-completion}
 \bigl(\E e^{p_1q_\pi(d+\mathsf S\zeta)}\bigr)^{1/p_1}
 =\det(I-M_\pi)^{-1/(2p_1)}
  \exp\sup_\zeta\left\{q_\pi(d+\mathsf S\zeta)
                                  -\frac{\|\zeta\|^2}{2p_1}\right\}.
\end{equation}
The supremum is bounded by the same $V_{L^{j_0}}(X,\phi)$ in
\eqref{eq:entry-observation-regulator}.  A broad energy has rank at most
its number of observation coordinates, as in
\cite[Equations~(83)--(84)]{F}.  There are at most
$C_{R_0}l(\pi)$ physical coordinates; proxy errors keep their small
coefficients.  Link-only cells enter through the reserved link cost,
not a broad rank.  The trace bound is consequently
\begin{equation}\label{eq:complex-entry-trace}
 \operatorname{tr}M_\pi\le C_{R_0}l(\pi)
       +C\sum_{\ell\in\pi}e^{-c(R_0+d_\ell)}
                                      \operatorname{poly}(d_\ell).
\end{equation}
A face contrast has uniformly bounded reference variance by the cell
flow bound, which proves the second term, even for arbitrarily long
links.  Since $\|M_\pi\|$ is bounded away from one,
$-\log\det(I-M_\pi)\le C\operatorname{tr}M_\pi$.
H\"older's inequality handles the ordinary regulator with its fixed
exponent reserve.  Translation by the bounded real part of $zf$ costs
$e^{C|X|}$.  Such a factor can be absorbed by weakening the gains of
\eqref{eq:entry-links}, because each link occupies only
$O(1+d_\ell/R_0)$ blocks.  The exponential gain per physical item is
not retained in this coarse estimate.

We have obtained, on the same field disks and in the same regulators,
\begin{equation}\label{eq:entry-complex-pattern}
 |F_\pi(b;\phi+zf)|
 \le e^{C_{R_0}l(\pi)}
       \prod_{\ell\in\pi}Ce^{-c(R_0+d_\ell)}
       \mathscr W_{L^{j_0}}^\kappa(X,\phi)e^{V_{L^{j_0}}(X,\phi)}.
\end{equation}
The domination also proves analyticity of each pattern after Gaussian
averaging.  All scalar denominators are local and nonzero, and the
integrals converge absolutely on compact parameter and field disks.

\paragraph{Recovering smallness in the parameter disk.}
For fixed $\pi,\phi,z,f$, apply the two-constants theorem to the upper
and lower half-disks $|b-\bc|<r_s$.  On the real diameter use
\eqref{eq:entry-real-pattern}; on the semicircle use
\eqref{eq:entry-complex-pattern}.  If $|b-\bc|\le\vartheta r_s$, the
harmonic measure of the semicircle is at most $\omega(\vartheta)$,
where $\omega(\vartheta)\to0$ as $\vartheta\downarrow0$.  The
regulator and a common weakened link product are independent of $b$
and can be divided out.  The resulting item factor is at most
\[
 \exp\{[-(1-\omega(\vartheta))cR_0
                         +\omega(\vartheta)C_{R_0}+C]l(\pi)\}.
\]
Choose $\vartheta>0$, depending only on $R_0$, so that its exponent is
at most $-cR_0l(\pi)/2$, after enlarging $R_0$ if needed.  Thus the
small pattern bound holds on the reduced complex parameter disk.
For patterns containing only links, both boundary estimates already
have the common link gains.

The fixed field disk has radius exceeding the polarization radius in
\eqref{eq:entry-point-norm}.  Cauchy's formula and polarization
therefore sum all field derivatives with that norm.  Connected
summation over the complete occupied covers, with the gains just
proved, is the same absolutely convergent sum as in
\cite[Theorem~4.6]{F}.  Its norm is at most $Ce^{-cR_0}$, uniformly in
$s$, the absolute entry index, and compatible volume.  The same bound
on compact parameter disks makes the sum Banach analytic.  Absorb
$\vartheta$ into $c_{R_0}$ and choose $R_0$ to reach the prescribed
$\rho$.  Parameter Cauchy estimates give \eqref{eq:entry-cauchy}.
The occupied covers retain every field dependency, as in the critical
construction.  The covariant placement and color rules, reference
Gaussian averaging, and field-negation identities are unchanged when
only the physical weights and centered scalar normalizers depend on
$b$.  Thus \eqref{eq:bulk-activity-symmetries} and the prescribed collar
rules hold throughout the parameter disk.  The exact finite algebra
and real positive normalizers give the partition representation for
real $b$, and the same local formulas give the copying identities.
\end{proof}

\subsection{The map at fixed reference}

The initialization puts the entire activity into the remainder.  At a
later side $L^j$, separate two singleton functions on each regular block:
\[
 e_B^0(\phi)=\frac12\sum_{v\in B}\sum_{e=\pm e_1,\pm e_2}
                                      |\nabla_e\phi(v)|^2,
 \qquad
 c_B^\alpha(\phi)=L^{-2j}\sum_{v\in B}\cos(\alpha\phi_v).
\]
Their coefficients are denoted by $t_j/2$ and $z_j$; the remaining
activity is $\mathcal R_j$.  Block covariance makes these coefficients
independent of the block.  Field negation excludes the sine singleton,
and square symmetry makes the retained neutral gradient quadratic
isotropic.  These are the two scalar directions in the bulk weak map
of \cite[Section~5.1]{F}.  The exact map integrates the shell
$\Gamma_{L^j}$, removes designated field-independent factors, and
reexpresses the result in the same subset partition algebra.  The
following interface collects the activity estimates used in the rest
of the paper.

\begin{proposition}[Activity map]\label{prop:map}
For all sufficiently large absolute indices $j$, the map on
$(t,z,\mathcal R)\in\mathbb C^2\times
\mathscr A^{\mathrm{bulk}}_{L^j,s}$, with the occupied-support and
designated-localization conventions of \cite[Proposition~4.1]{F},
preserves the bulk symmetry class and is analytic on a fixed ball
$|t|+|z|+\|\mathcal R\|_{L^j,s}<\rho_{\mathrm{map}}$.
Its linear part is
\begin{equation}\label{eq:map-triangular}
 (t,z,\mathcal R)\longmapsto
   (t+P_j^t\mathcal R,\lambda_jz+P_j^z\mathcal R,
                                                \mathsf C_j\mathcal R),
 \quad \|\mathsf C_j\|\le\theta_0<1,
 \quad \lambda_j=1+O(e^{-cj}),
\end{equation}
with uniformly bounded rows $P_j^t,P_j^z$.  The nonlinear part and its
Taylor derivatives have quadratic bounds with uniform constants.

There are analytic normal coordinates $(u,Z,Q)$, uniform in the entry
side, such that the pure Gaussian curve is $(u,0,0)$ and the map fixes
$u$ on that curve.  With $b_*=2\log L$ and a constant $H>0$, define
\begin{equation}\label{eq:balanced-coordinates}
                  x=b_*u,\qquad y=\sqrt{b_*H}\,Z,\qquad W=(x,y).
\end{equation}
The coordinate changes obey
\begin{equation}\label{eq:coordinate-comparison}
 t=u+O(u^2+\|Q\|),\qquad
 z=(1+O(e^{-cj}))Z+O(u^2+\|Q\|),\qquad
 \mathcal R=O(u^2)+Q,
\end{equation}
with the differentiated analytic bounds of these orders.  In a fixed
small tube,
\begin{equation}\label{eq:normal-map}
 \begin{split}
 W_+&=W+(-y^2,-xy)+\mathcal E_j(W,Q),\\
 |\mathcal E_j(W,Q)|
 &\le C\{j^{-2}|W|^2+|W|^3+(|W|+\|Q\|)\|Q\|\},\\
 \|Q_+\|&\le\theta\|Q\|+C|W|^2,\qquad \theta<1.
 \end{split}
\end{equation}
The error is an analytic function with the corresponding Taylor
bounds on all fixed derivatives.

On a separately initialized compatible torus with a fixed range buffer,
the two singleton coefficients copy the plane coefficients through its
stop.  The torus remainder has its own contraction and analytic
quadratic bound.  Its derivative at zero has no singleton forcing.
All constants are uniform in the stop, the entry side, and their
associated norms.
\end{proposition}

\begin{proof}
The exact activity map and its uniform norm bounds are
\cite[Proposition~4.1]{F}; the triangular form is
\cite[Equation~(93)]{F}.  We recall the coordinate construction to
specify the derivative bounds being used.  Put
$\ell_j=\prod_{h\ge j}\lambda_h=1+O(e^{-cj})$ and
$\mathsf C_{h:j}=\mathsf C_{h-1}\cdots\mathsf C_j$.  The balanced
linear coordinates are
\[
 T=t+\sum_{h\ge j}P_h^t\mathsf C_{h:j}\mathcal R,
 \qquad
 Z=\ell_jz+\sum_{h\ge j}\ell_{h+1}P_h^z\mathsf C_{h:j}\mathcal R.
\]
The series converge geometrically.  At the chosen entry, initialize an
exact Gaussian interaction $\exp(d\sum_Be_B^0/2)$ and iterate the
same map.  Its balanced gradient is $T_j^g(d)=d+O(d^2)$ and its
remainder is $\mathcal R_j^g(d)=O(d^2)$, uniformly on a common complex
disk.  Define $u=(T_j^g)^{-1}(T)$ and
$Q=\mathcal R-\mathcal R_j^g(u)$, as in
\cite[Equations~(95), (99), and (101)]{F}.  This gives
\eqref{eq:coordinate-comparison}; in particular its linear remainder
rows are bounded, while derivatives of its quadratic scalar terms are
$O(|u|)$.

The parity estimates of \cite[Lemma~5.2]{F} give leading scalar terms
$-H_jZ^2$ and $-b_juZ$, with no nonconstant term on the Gaussian curve.
Their remainders have the differentiated orders used in
\eqref{eq:normal-map}.  Moreover
\cite[Lemma~5.3]{F} gives
$|H_j-H|+|b_j-b_*|\le C(1+j)^{-2}$.
Rescaling by \eqref{eq:balanced-coordinates} yields the displayed
map.  Terms with a small coefficient multiplying $Q$ are absorbed into
a slightly larger contraction constant.  The copying assertion and the
separate torus contraction are part of the same local-map statement;
no equality between plane and torus remainders is required.
\end{proof}

The triangular contraction also has a useful elementary consequence.
In equivalent scalar and remainder norms, the derivative of the full
map on a sufficiently small tube has norm at most $1+\gamma$, for any
prescribed $\gamma>0$.  To see this, first use the balanced coordinates,
whose scalar linear rows contain no remainder, and choose a product
norm with the fixed remainder contraction.  At zero the operator norm
is at most one.  Uniform analytic derivative bounds then give the
claim on a small tube.  A torus remainder driven by these scalar
coordinates satisfies the corresponding growth bound by convolution.

\subsection{Finite torus observables and comparison of observation grids}

Fix a sufficiently large dyadic range buffer $D_0$ and, for an absolute
stop $j$, use the torus of side
\begin{equation}\label{eq:stopping-tori}
                              n_j=D_0L^j.
\end{equation}
Each such torus is initialized separately at side $qR_0s$.  It has
$D_0^2$ blocks at its stop and its singletons copy the plane history.
The physical cutoff initialization requires $n_j\ge sP_0$.
Define the physical first cosine and its noisy-cell approximation by
\begin{equation}\label{eq:physical-cosine}
 \begin{split}
 H_n(h)&=n^{-2}\sum_{v\in(\Z/n\Z)^2}h_v\cos(2\pi v_1/n),\\
 X_{n,s}&=\langle\eta,f_{n,s}\rangle,
 \qquad f_{n,s}(Q)=(s/n)^2\cos(2\pi x_{Q,1}/n),
 \end{split}
\end{equation}
where $x_Q$ is the arithmetic center of the primary sites in $Q$.
The origins of all grids agree.  Both real tests annihilate constants.
The phase test is $e^{i\bar\eta}$, where $\bar\eta$ is the uniform
mean of the cell observations, interpreted modulo $2\pi$.
Set $\mathfrak g=1/(8\pi^2)$ and let
\[
 v^0_{n,s}=\langle f_{n,s},(I+aB_sA^+B_s^*)f_{n,s}\rangle
\]
be the observation variance in the fixed Gaussian reference.

\begin{proposition}[Analytic torus observables]\label{prop:readouts}
On a fixed small complex ball of terminal activities, the cutoff
variance correction and the cutoff phase expectation extend to analytic
functions of the activity coordinates.  If
$\nu_j=|t_j|+|z_j|+\|\mathcal R^{\rm tor}_{j|j}\|$, they satisfy
\begin{equation}\label{eq:torus-readouts}
 \begin{split}
 \widehat\Var(X_{n_j,s})-v^0_{n_j,s}
      &=\kappa_{n_j,s}t_j+O(\|\mathcal R^{\rm tor}_{j|j}\|+\nu_j^2),\\
 \widehat\E e^{i\bar\eta}
      &=p_{j,s}z_j+O(\|\mathcal R^{\rm tor}_{j|j}\|+\nu_j^2).
 \end{split}
\end{equation}
All fixed activity derivatives have the corresponding Taylor bounds,
uniformly in entry and stop.  The multipliers are positive, bounded
above and below, and
\begin{equation}\label{eq:readout-multipliers}
 v^0_{n,s}=a\mathfrak g+O(s/n+n^{-2}),\qquad
 \kappa_{n,s}=a\mathfrak g+O(s/n+n^{-2}).
\end{equation}
At a common absolute stop, the multipliers $p_{j,s}$ for different
entries agree except for their mean-observation-noise factor; their
relative difference is $O((s_1/n)^2+(s_2/n)^2)$.
\end{proposition}

\begin{proof}
These are the terminal analytic functionals of
\cite[Equations~(97)--(98), (116)--(119), and (130)]{F}.
Their analyticity does not require a scaling limit for the physical
input.  Indeed, on the fixed $D_0^2$-block grid the Gaussian and common-mean
integral of the terminal subset sum is $1+O(D_0^2\nu_j)$, uniformly
on a fixed complex source disk.
Choose the activity ball so that this error is less than $1/2$.
Both normalized numerator and denominator are then analytic and the
denominator is nonzero.

Completing the Gaussian source $wX_{n,s}$ before entry averaging gives
the factor $e^{v^0_{n,s}w^2/2}$ and translates the interaction by a
microscopic cosine $F(w)=w\gamma_{n,s}\cos(2\pi v_1/n)$.
Writing $\beta_{n,s}$ for the cell-average cosine factor,
\[
 B_s\cos(2\pi v_1/n)=\beta_{n,s}\cos(2\pi x_{Q,1}/n),\qquad
 \gamma_{n,s}=\frac{2v^0_{n,s}}{\alpha\beta_{n,s}},\qquad
 \beta_{n,s}=1+O((s/n)^2).
\]
The linear gradient singleton contributes
$t_j\langle F,AF\rangle/2$ to the logarithm; mean integration removes
the first harmonic at linear order.  Since
$\langle\cos(2\pi v_1/n),A\cos(2\pi v_1/n)\rangle
=2n^2\sin^2(\pi/n)$, this gives the first line and
\eqref{eq:readout-multipliers}.  The remainder costs its norm and all
higher terms cost $O(\nu_j^2)$.

For the phase insertion the common-mean integral instead selects the
fundamental singleton.  Its coefficient is
\[
 p_{j,s}=\frac{D_0^2}{2}
           \exp\{\tau_{j,s}/2-aD_{{\rm tail},j}/2\}>0,
\]
where $D_{{\rm tail},j}$ is the diagonal of the retained mean-free
covariance and $\tau_{j,s}$ is the variance, in height units, of the
integrated high mean including the observation noise.  Both are bounded
on the fixed terminal geometry.  At the same absolute stop the high
covariances telescope to the same constant-mode continuation.  The only
entry-dependent contribution to $\tau_{j,s}$ is the noise variance
$(s/n)^2$.  This proves the second assertion and the comparison of
multipliers.  Cauchy estimates on a smaller activity ball give all
stated derivatives.
\end{proof}

The next two lemmas permit comparison through these observables without
identifying the activity spaces of distinct observation grids.

\begin{lemma}[Cutoff error and its derivative]\label{lem:cutoff-derivative}
For real $b$ in the interval of Proposition~\ref{prop:entry}, and
$n\ge sP_0$ with $\log n=o(s^2)$, replacing the original torus
expectation by its normalized cutoff expectation changes the variance
of $X_{n,s}$ and the phase expectation by at most
\begin{equation}\label{eq:cutoff-derivative-error}
                                Cn^Ce^{-cs^2}.
\end{equation}
The same bound holds for their first $b$ derivatives and for divided
differences between any two real parameters in that interval.
\end{lemma}

\begin{proof}
The cell-flow proof of \eqref{eq:entry-large-field} also bounds the
probability that a prescribed face difference exceeds $s/2$ by
$Ce^{-cs^2}$, uniformly in the compact $b$ interval.  On its complement,
$1-\chi_s$ has that bound.  Summing over faces and applying H\"older
with the uniform reference moments gives the unnormalized loss for the
phase and the first two moments of $X_{n,s}$, as in
\cite[Lemma~4.4]{F}.  The cutoff mass is consequently $1-o(1)$ after
normalization by the uncut law.

To differentiate a normalized expectation, insert the Poisson score
$N/b$.  Lemma~\ref{lem:factorial} bounds every fixed moment of $N$ by a
fixed power of $n$, uniformly on the parameter interval.  Apply
H\"older to the product of this insertion, the observable, and the
cutoff loss.  The face tail retains a bound $e^{-c's^2}$, while the
other factors cost a power of $n$.  Differentiating the cutoff
normalizer has the same estimate because it is bounded away from zero.
This proves \eqref{eq:cutoff-derivative-error} for derivatives.
Integrating the derivative bound proves the divided-difference version.
\end{proof}

\begin{lemma}[Observation-grid comparison]\label{lem:grid-comparison}
After dividing the phase expectation by its deterministic mean-noise
factor, it is exactly $\E e^{i\bar h}$, independently of $s$.
After subtracting the independent observation-noise variance and dividing
by $\beta_{n,s}^2$, the variance of $X_{n,s}$ differs from
$\Var H_n$ by $O((s/n)^4)$.  Its first $b$ derivative, and its divided
differences over a compact real parameter interval, differ by at most
\begin{equation}\label{eq:grid-derivative-error}
                                  Cn^2(s/n)^4.
\end{equation}
The estimates are uniform for compatible grids with $n/s$ sufficiently
large.  Cutoff observables incur the additional error in
\eqref{eq:cutoff-derivative-error}.
\end{lemma}

\begin{proof}
The spatial mean of $B_sh$ is exactly the microscopic mean.  The
independent mean noise has variance $(s/n)^2$, proving the phase
identity.  The profile $B_s^*f_{n,s}$ projects onto the microscopic
first cosine with coefficient $\beta_{n,s}$ and has no first sine
component, by reflection around arithmetic cell centers.  The remaining
Fourier modes have frequency-index distance at least $cn/s$ from zero.
Their total squared coefficient norm is at most
$Cn^{-2}(s/n)^2$, whereas $A^+$ restricted to these modes has norm at
most $Cs^2$.  Their reference cost is therefore $O((s/n)^4)$.
Translation invariance of the physical torus law makes distinct Fourier
modes exactly orthogonal.  Thus the variance is the first-cosine
variance plus the alias variance, without a cross term.  The reference
bound controls the latter variance.  Lemma~\ref{lem:response}, applied
to the alias test, controls its derivative by its reference cost times
$|E|=2n^2$.  The factor $\beta_{n,s}^{-2}$ is uniformly bounded and
independent of $b$.  Integration gives the divided-difference estimate.
\end{proof}

\subsection{The critical histories and late entry}

We finish the interface by recording exactly which critical information
is used later.  A history means the plane sequence initialized wholly
in its remainder and iterated with Proposition~\ref{prop:map}, together
with the separately initialized tori at its valid stops.

\begin{proposition}[Critical histories]\label{prop:critical}
The fixed parameters and entry threshold can be chosen so that there
is a family $\mathcal H^{(k)}$ at $b=\bc$, with admissible
$s_k\asymp2^k$ and entry indices $j_0(k)=O(k)$, having the following
properties.  For every fixed $0<c<c'<3$, uniformly on
$c2^k\le j\le c'2^k$,
\begin{equation}\label{eq:critical-asymptotic}
 W_j^{(k)}=\frac{(1,1)}j+O\!\left(\frac{\log(2+j)}{j^2}\right),
 \qquad
 \|Q_j^{(k)}\|+\|\mathcal R_j^{(k)}\|=O(j^{-2}).
\end{equation}
The separately initialized torus remainders also satisfy $O(j^{-2})$
on these stage ranges behind every valid later stop in the same
history up to $4\cdot2^k$.

More generally, for any prescribed sufficiently small tube and any
integer bound $J_{\max}(s)=o(s^2)$, the critical plane history at an
admissible side $s$ stays in that tube through $J_{\max}(s)$, for all
large $s$.  Every stage of its separately initialized tori with valid
stops below $J_{\max}(s)$ stays in the corresponding small ball.  The
choices of geometric parameters are independent of $s$ and
$J_{\max}$.
\end{proposition}

\begin{proof}
The family and its confinement through $4\cdot2^k$ are
\cite[Proposition~5.1 and Lemma~5.5]{F}.  We explain both the
second-coordinate precision and the extension to longer critical
histories.

In the proof of that proposition, let $Y_j=|Z_j|$ for the spliced
critical histories.  Equations~(135)--(137) of \cite{F} give
\[
 u_j\asymp Y_j\asymp j^{-1},\qquad
 u_j=\sqrt{H/b_*}\,Y_j+O(j^{-2}),\qquad
 u_j=(b_*j)^{-1}+O(\log(2+j)/j^2).
\]
The overlap bounds in \cite[Lemma~5.7]{F} transfer both coordinates
between histories with $O(j^{-2})$ error on the indicated fractional
ranges.  The phase formula \cite[Equation~(130)]{F}, together with
positive physical characteristic coefficients and the exponentially
small cutoff loss, selects $Z_j>0$ there.  Rescaling by
\eqref{eq:balanced-coordinates} gives \eqref{eq:critical-asymptotic}.
The plane remainder bound is part of the same proposition and its odd
cone estimate.  For a larger torus, its separate contraction is driven
by the copied scalar sequence.  Contracting from an earlier fractional
range suppresses its initial contribution exponentially and sums an
$O(j^{-2})$ forcing, giving the asserted bound at every intermediate
stage.

For the last assertion, the physical input to critical confinement is
\cite[Lemma~5.4]{F}.  It gives the critical torus Gaussian first-cosine
limit and factorization with the mean phase.  For any fixed tolerance,
first $P_0$ and then $s$ can be chosen so that the variance and mixed
phase discrepancies are below that tolerance for every compatible
$n\ge sP_0$.  No rate of convergence in $n$ is asserted or needed.
The mixed phase functional has a nonzero multiplier on $z$, as shown
in \cite[Equations~(118)--(119)]{F}.  The variance multiplier in
Proposition~\ref{prop:readouts} is nonzero as well.

Suppose there were a first exit from a tube of radius $\rho$.
Before that exit both plane and torus remainder contractions give
\[
 \|\mathcal R_j\|+\|\mathcal R^{\rm tor}_{j|j}\|
                    \le C\theta^{j-j_0}\rho_{\rm ent}+C\rho^2,
\]
where $\rho_{\rm ent}$ is the entry norm.  At the candidate exit the
bounded map leaves all states inside a fixed multiple of the tube, so
the terminal observables are still defined.  The variance and mixed
phase identities then bound the two singleton coefficients by
$C(\rho_{\rm ent}+\rho^2+\delta_{\rm phys}+\delta_{\rm cut})$.
Here $\delta_{\rm phys}$ is the uniform critical discrepancy just
described.  The number of initial indices with $n_j<sP_0$ is at most
$\max(0,\lceil\log_L(R_0^4/(D_0q))\rceil)$, independent of $s$;
their amplification is absorbed by the small entry norm.
Finally, uniformly for $j\le J_{\max}(s)$,
\[
 \delta_{\rm cut}\le Cn_j^Ce^{-cs^2}
           \le \exp\{O(J_{\max}(s))-cs^2\}\longrightarrow0.
\]
Choose $\rho$ small, then the entry and physical discrepancies small
relative to $\rho$, and finally $s$ large.  The preceding bound
excludes the first exit.  Separate torus contractions give the same
conclusion at all earlier stages of each valid stopped torus.  This is
the finite-history argument of \cite[Lemma~5.5]{F} with precisely the
larger range permitted by the cutoff estimate.
\end{proof}
